% GENERATED FILE. Edit canonical lessons, then run npm run generate:books.
% canonical-source-hash: fnv1a64:be5b49c4b93d8735
% canonical-lessons: GU-C477-odkar, GU-C477-aphro, GU-C477-ubko, GU-C477-okvu, GU-C477-jiravvu

\chapter{After a Big Meal}
\label{ch:gu-after-a-big-meal}
\hlchaptermodality{477}

\begin{chapteropening}
\textbf{By the end of this chapter:} \emph{I can say what my stomach feels after eating.}

Complete the last lesson of chapter 477: I can say what my stomach feels after eating.
\end{chapteropening}

\section[\texorpdfstring{\gu{ઓડકાર} (oḍkār)}{ઓડકાર (oḍkār)}]{\gu{ઓડકાર} (oḍkār) — a belch}

\label{lesson:GU-C477-odkar}



\begin{quote}
Before the new one: say the Gujarati for to lie down to sleep, then the Gujarati for sleeplessness.
\end{quote}

\subsection*{You'll want to know: \gu{ઓડકાર}}
\textbf{\gu{ઓડકાર}} — \emph{oḍkār} — \textquotedblleft{}a belch\textquotedblright{}.

\begin{quote}
\gu{જમ્યા} \gu{પછી} \gu{દાદાએ} \gu{મોટો} \gu{ઓડકાર} \gu{ખાધો}.
\end{quote}

\begin{quote}
\emph{jamyā pachī dādāe moṭo oḍkār khādho.}
\end{quote}

\begin{quote}
After the meal my grandfather gave a loud belch.
\end{quote}

A masculine noun.

\begin{cousinweb}
\textbf{\gu{ઓડકાર}} (\emph{oḍkār}) probably goes back to Sanskrit \emph{udgāra}, \textquotedblleft{}the act of discharging, spitting out\textquotedblright{}, as Monier-Williams records it.
\end{cousinweb}

\begin{grammarlens}[title={What you've built}]
One more word for describing what you lived and felt.
\end{grammarlens}

\subsection*{Your turn}
\begin{itemize}
\raggedright
  \item \emph{Say it:} \emph{oḍkār}
  \item \emph{Say it:} \emph{oḍkār}, once more
  \item \emph{Say it:} \emph{jamyā pachī dādāe moṭo oḍkār khādho.} — After the meal my grandfather gave a loud belch.
  \item \emph{Recall:} say the Gujarati for to lie down to sleep, then the Gujarati for sleeplessness, then say \emph{oḍkār} again
\end{itemize}

\begin{tcolorbox}[breakable,colback=teal!4,colframe=teal!35!black,title={Before you move on}]
What does \gu{ઓડકાર} mean? (\textquotedblleft{}A belch\textquotedblright{}.) Say it once more.
\end{tcolorbox}
\section[\texorpdfstring{\gu{આફરો} (āphro)}{આફરો (āphro)}]{\gu{આફરો} (āphro) — bloating}

\label{lesson:GU-C477-aphro}



\begin{quote}
Before the new one: say the Gujarati for sleeplessness, then the Gujarati for a belch.
\end{quote}

\subsection*{You'll want to know: \gu{આફરો}}
\textbf{\gu{આફરો}} — \emph{āphro} — \textquotedblleft{}bloating\textquotedblright{}.

\begin{quote}
\gu{વધારે} \gu{ભોજન} \gu{પછી} \gu{પેટમાં} \gu{આફરો} \gu{ચઢ્યો}.
\end{quote}

\begin{quote}
\emph{vadhāre bhojan pachī peṭmā̃ āphro caḍhyo.}
\end{quote}

\begin{quote}
After the big meal my stomach felt bloated.
\end{quote}

A masculine noun.

\begin{cousinweb}
\textbf{\gu{આફરો}} (\emph{āphro}) goes back to a reconstructed \emph{*āsphara-}, \textquotedblleft{}swelling\textquotedblright{}; CDIAL lists the Gujarati word, \textquotedblleft{}wind in belly\textquotedblright{} (CDIAL 1526).
\end{cousinweb}

\begin{grammarlens}[title={What you've built}]
One more word for describing what you lived and felt.
\end{grammarlens}

\subsection*{Your turn}
\begin{itemize}
\raggedright
  \item \emph{Say it:} \emph{āphro}
  \item \emph{Say it:} \emph{āphro}, once more
  \item \emph{Say it:} \emph{vadhāre bhojan pachī peṭmā̃ āphro caḍhyo.} — After the big meal my stomach felt bloated.
  \item \emph{Recall:} say the Gujarati for sleeplessness, then the Gujarati for a belch, then say \emph{āphro} again
\end{itemize}

\begin{tcolorbox}[breakable,colback=teal!4,colframe=teal!35!black,title={Before you move on}]
What does \gu{આફરો} mean? (\textquotedblleft{}Bloating\textquotedblright{}.) Say it once more.
\end{tcolorbox}
\section[\texorpdfstring{\gu{ઉબકો} (ubko)}{ઉબકો (ubko)}]{\gu{ઉબકો} (ubko) — a retch; a wave of nausea}

\label{lesson:GU-C477-ubko}



\begin{quote}
Before the new one: say the Gujarati for a belch, then the Gujarati for bloating.
\end{quote}

\subsection*{You'll want to know: \gu{ઉબકો}}
\textbf{\gu{ઉબકો}} — \emph{ubko} — \textquotedblleft{}a retch; a wave of nausea\textquotedblright{}.

\begin{quote}
\gu{બસમાં} \gu{મીરાને} \gu{ઉબકા} \gu{આવ્યા}.
\end{quote}

\begin{quote}
\emph{basmā̃ Mīrāne ubkā āvyā.}
\end{quote}

\begin{quote}
Mira felt sick on the bus.
\end{quote}

A masculine noun.

\begin{cousinweb}
\textbf{\gu{ઉબકો}} (\emph{ubko}) goes back to a reconstructed \emph{*ubbakka-}, \textquotedblleft{}vomited\textquotedblright{}; CDIAL lists the Gujarati word, \textquotedblleft{}retching\textquotedblright{} (CDIAL 2337).
\end{cousinweb}

\begin{grammarlens}[title={What you've built}]
One more word for describing what you lived and felt.
\end{grammarlens}

\subsection*{Your turn}
\begin{itemize}
\raggedright
  \item \emph{Say it:} \emph{ubko}
  \item \emph{Say it:} \emph{ubko}, once more
  \item \emph{Say it:} \emph{basmā̃ Mīrāne ubkā āvyā.} — Mira felt sick on the bus.
  \item \emph{Recall:} say the Gujarati for a belch, then the Gujarati for bloating, then say \emph{ubko} again
\end{itemize}

\begin{tcolorbox}[breakable,colback=teal!4,colframe=teal!35!black,title={Before you move on}]
What does \gu{ઉબકો} mean? (\textquotedblleft{}A retch; a wave of nausea\textquotedblright{}.) Say it once more.
\end{tcolorbox}
\section[\texorpdfstring{\gu{ઓકવું} (okvũ)}{ઓકવું (okvũ)}]{\gu{ઓકવું} (okvũ) — to throw up}

\label{lesson:GU-C477-okvu}



\begin{quote}
Before the new one: say the Gujarati for bloating, then the Gujarati for nausea.
\end{quote}

\subsection*{You'll want to know: \gu{ઓકવું}}
\textbf{\gu{ઓકવું}} — \emph{okvũ} — \textquotedblleft{}to throw up\textquotedblright{}.

\begin{quote}
\gu{બાળકે} \gu{દૂધ} \gu{ઓકી} \gu{નાખ્યું}.
\end{quote}

\begin{quote}
\emph{bāḷke dūdh okī nākhyũ.}
\end{quote}

\begin{quote}
The baby threw up its milk.
\end{quote}

Past, with \gu{તેણે} (he or she …): \gu{તેણે} \gu{દૂધ} \gu{ઓક્યું}, \emph{teṇe dūdh okyũ}.

\begin{cousinweb}
\textbf{\gu{ઓકવું}} (\emph{okvũ}) goes back to a reconstructed \emph{*ōkk-}, \textquotedblleft{}vomit\textquotedblright{}, an imitative word; CDIAL lists the Gujarati verb (CDIAL 2538).
\end{cousinweb}

\begin{grammarlens}[title={What you've built}]
One more word for describing what you lived and felt.
\end{grammarlens}

\subsection*{Your turn}
\begin{itemize}
\raggedright
  \item \emph{Say it:} \emph{okvũ}
  \item \emph{Say it:} \emph{okvũ}, once more
  \item \emph{Say it:} \emph{bāḷke dūdh okī nākhyũ.} — The baby threw up its milk.
  \item \emph{Recall:} say the Gujarati for bloating, then the Gujarati for nausea, then say \emph{okvũ} again
\end{itemize}

\begin{tcolorbox}[breakable,colback=teal!4,colframe=teal!35!black,title={Before you move on}]
What does \gu{ઓકવું} mean? (\textquotedblleft{}To throw up\textquotedblright{}.) Say it once more.
\end{tcolorbox}
\section[\texorpdfstring{\gu{જીરવવું} (jīravvũ)}{જીરવવું (jīravvũ)}]{\gu{જીરવવું} (jīravvũ) — to stomach; to digest (food or bad news)}

\label{lesson:GU-C477-jiravvu}



\begin{quote}
Before the new one: say the Gujarati for nausea, then the Gujarati for to throw up.
\end{quote}

\subsection*{You'll want to know: \gu{જીરવવું}}
\textbf{\gu{જીરવવું}} — \emph{jīravvũ} — \textquotedblleft{}to stomach; to digest (food or bad news)\textquotedblright{}.

\begin{quote}
\gu{આવું} \gu{દુઃખ} \gu{જીરવવું} \gu{અઘરું} \gu{છે}.
\end{quote}

\begin{quote}
\emph{āvũ duḥkh jīravvũ agharũ che.}
\end{quote}

\begin{quote}
Such sorrow is hard to bear.
\end{quote}

Past, with \gu{તેણે} (he or she …): \gu{તેણે} \gu{દુઃખ} \gu{જીરવ્યું}, \emph{teṇe duḥkh jīravyũ}.

\begin{cousinweb}
\textbf{\gu{જીરવવું}} (\emph{jīravvũ}) goes back to a reconstructed \emph{*jīrayati}, \textquotedblleft{}digests\textquotedblright{}; CDIAL lists the Gujarati verb, \textquotedblleft{}to digest\textquotedblright{} (CDIAL 5236).
\end{cousinweb}

\begin{grammarlens}[title={What you've built}]
One more word for describing what you lived and felt.
\end{grammarlens}

\subsection*{Your turn}
\begin{itemize}
\raggedright
  \item \emph{Say it:} \emph{jīravvũ}
  \item \emph{Say it:} \emph{jīravvũ}, once more
  \item \emph{Say it:} \emph{āvũ duḥkh jīravvũ agharũ che.} — Such sorrow is hard to bear.
  \item \emph{Recall:} say the Gujarati for nausea, then the Gujarati for to throw up, then say \emph{jīravvũ} again
\end{itemize}

\begin{tcolorbox}[breakable,colback=teal!4,colframe=teal!35!black,title={Before you move on}]
What does \gu{જીરવવું} mean? (\textquotedblleft{}To stomach; to digest (food or bad news)\textquotedblright{}.) Say it once more.
\end{tcolorbox}
